\section{Three Types of Ascesis}

In an accordance with the hierarchy spiritus, anima, corpus (spirit, soul, body), we distinguish three types of ascesis, three types of self-control: ascesis of spirit, ascesis of soul, ascesis of body. Spiritual ascesis or mental ascesis is an ascesis of thoughts, a mental self-control. Ascesis of a soul represents control over a character. [1] Bodily ascesis means control of an action, conduct and behaviour.

A key to successful practice of an ascesis is to practice in an accordance with the given hierarchy: Every impulse attacks first one's mind in form of a thought. A bad thought leads to a character disharmony, e.g., desire. Yet the desire causes wrong acting (and not the other way). This is why the beginning of every successful ascesis is to suppress a wrong thought, a bad mental impulse. [2]

In order to suppress incoming thoughts, much advice is given by saints and holy teachers: do not believe your thoughts, ignore them, do not pay them attention etc. In this manner one purifies his mind and as a by-effect deepens his consciousness.

If one does not respect the hierarchy and lets into their mind all kinds of forms, shapes, images and thoughts, and yet subsequently tries to fight them and resist, they become nothing more than overlaid, overshadowed or displaced into one's unconsciousness, which leads to a dangerous polarity that will sooner or later burst out again from within. [3] If one is not aware of this and engages himself in a never-ending fight with the constant inflow of shapes, images and temptations, one only weakens and exhausts himself. A loss is thus only a question of time.

To stop identifying with thoughts and to annihilate all the wrong thoughts immediately in the very beginning is the key to every ascetic and thus to a victory in the greatest holy war[4], where the reward is the awakening and approaching to the centre of the world, which is a residence to the great peace.


\hfill

\textbf{Notes:}

[1] Some authors not arbitrarily call the soul the emotional body or astral body (aster = star).

[2] This is why there is no difference between a sinful thought, a desire for a sin, and a sinful act. Men are constantly creating their Destiny through the quality and power of their thoughts: great desire causes great obstacles to arise. Purify your mind, thus straighten your path.

[3] This affects mostly the beginners who cannot yet see the crucial difference between darkness and emptiness.

[4] Al-Jihad al-akbar from an Islamic hadith means literally the greatest holy war. This is also one of the outer meanings of Jesus' walking on water: the immediate domination over the world of shapes and forms. This superiority is an implicit condition to every initiation. See also Ephesians 6:11-19.



\flrightit{Posted on 2013-11-11 by Mars Ultor }

\begin{center}* * *\end{center}

\begin{footnotesize}\begin{sffamily}



\texttt{X on 2013-11-12 at 01:11 said: }

“Ascesis of a soul represents control over a character.”

Can you please explain this detailedly?

“This affects mostly the beginners who cannot yet see the crucial difference between darkness and emptiness.”

What is the exact difference between darkness and emptiness? I am a beginner.


\hfill

\texttt{Kidd on 2013-11-12 at 12:49 said: }

Where does the will fit in? It is something that acts upon either of the three domains? Also is it not possible for proper thoughts and desires to follow from proper actions/forms (aka the physical sphere)?


\hfill

\texttt{paulo on 2013-11-12 at 13:30 said: }

And how about replace bad thoughts with good thoughts?Is this not part of the process?If some one is a misanthrope and hate people,he is constantly thinking about violence.He should replace with images off peace?


\hfill

\texttt{Mars on 2013-11-13 at 03:31 said: }

I know the article is quite brief. Its intention is to show the key to an ascesis. It is meant as an advice for those who practice it. Now, to answer the questions:

Control over a character means to force yourselves to keep only positive affections, and thus to develop positive personal traits and character. Yes, as Kidd says, character and deeds retroactively influence one's thoughts, but the initial impulse is always mental.

Darkness X emptiness: Observe your thoughts very slowly, gently and carefully until you start sensing the main source of every thought. Gradually your intuition (in-tueri = insight) intensifies and you will find out easily.

Kidd: Yes, the will is a means to every ascesis – to suppress the wrong thoughts. But it is quite impossible to achieve emptiness through will.

Paolo: Yes, it is a part of the process. He is not thinking about violence because he is a misanthrope. He is a misanthrope because he keeps the violent fantasies.


\hfill

\texttt{Pickman on 2013-11-14 at 07:06 said: }

I love how these generated avatars capture with uncanny semblance the character of each poster (even my own). 

I have an open question for you all and any observers reading this: Has anyone here actually mastered these thought control techniques yet? 

I extend that question to Cologero in particular. In this modern age of turbulence and fast living how are we to implement these or what would be the ideal way. In other words what would it take if one had the absolute will and inclination to achieve these occult states?


\hfill

\texttt{Pickman on 2013-11-14 at 07:14 said: }

May I also recommend a list of authors in the Western tradition who adhered to such techniques from ancient to modern (there is a long list) in comparison to the Eastern doctrines of a similar nature. Self-mastery and it's practical application is imperative for us.

This series of articles on Evola's purification metaphysics has been most inspiring or most re-minding and more of the same would be welcomed.


\hfill

\texttt{scardanelli on 2013-11-14 at 10:37 said: }

I don't necessarily see why the task or its implementation would necessarily be any different due to the times we are living in. The essential nature of things has not changed…we are still ruled by our thoughts, appetites, sex… begin by beginning. 

The thought that we need special consideration before applying these techniques is just a barrier to beginning. One can fret about what brand, style, and level of arch support of his running shoes before training for a marathon, or you can just start running.


\hfill

\texttt{Pickman on 2013-11-14 at 15:14 said: }

Usually in these “explanations” we have a brief overview of the exercise followed by a theoretical explanation of its effect on the spirit. However it is rarely detailed how to actually achieve results. 

Mars how about more specific details to the beginner with an overview to overcoming the possible pitfalls that occur.

Scardanelli: My point is what environment is most preferred for optimal results as a primer. Yet lets just assume that one is active within the world, then an indestructible interiority must be established while the outward persona persists in mundane matters. This is difficult to accomplish, yet it echos what what Evolas autarky (and other mystic's i.e, Eckhart) was all about – “living in the world yet not being a part of it”.


\hfill

\texttt{Kid on 2013-11-15 at 16:04 said: }

Well this era would be more difficult because of the amount and the nature of the stimulus in our environment. I mean, even ascetics of eons past felt the urge to become forest dwellers to be away from it all! Imagine how they'd feel about today's world!


\hfill

\texttt{Cologero on 2013-11-17 at 19:44 said: }

Pickman, no, I cannot walk through fire unscathed. Nevertheless, no effort is wasted for one with an absolute will and inclination. You strive to attain whatever states are justly yours. Plotinus, for example, attained the ecstatic union with the One four times in his life. Was that worth the effort?

The modern world, contrary to some comments expressed here, provides ample opportunities for spiritual development. The modern world, as such, is unreal and is based solely on appearances. In our time, we can actually see and experience what was previously given as prophecies centuries and millennia ago. Even the social critiques from 100 years ago have come to pass in a concrete and unmistakable way in our time. Once you can dissipate the hypnotic fog of the modern world, many things will be revealed to you.

The Owl of Minerva flies at dusk, so in a sense we can become wiser because we can see the whole cycle unfolding. Those who came earlier may have had it too easy and never made the correct efforts. How do you suppose we came to this low point? Certainly not because our predecessors were keeping tradition alive.

We do know that this age will end and a new world will begin. Certain old souls are among us to prepare the way. Do you remember who you are and what you need to accomplish? Why do you bother to read some obscure blog like Gornahoor? Answer those questions and many things will be clear.


\hfill

\texttt{Logres on 2013-11-18 at 17:08 said: }

“This affects mostly the beginners who cannot yet see the crucial difference between darkness and emptiness.”

This comment was very, very helpful \& timely for me.

In fact the whole article has been helpful – thanks, Mars Ultor. You may have saved me a great deal of trouble \& time.


\end{sffamily}\end{footnotesize}
